\documentclass[a4paper, 12pt]{article}
\usepackage{fontspec}
\usepackage{xeCJK}
\usepackage{amsmath}
\setCJKmainfont{Noto Sans CJK SC}

\usepackage{hyperref}
\usepackage{algorithm}
\usepackage[noend]{algpseudocode}
\usepackage{mdframed}
\usepackage{xcolor}
\usepackage{amsthm}
\usepackage{graphicx}
\usepackage{booktabs}
\usepackage{enumitem}
\usepackage[most]{tcolorbox}
\usepackage[protrusion=true,expansion=false]{microtype}
\emergencystretch=2em

\title{Interactive Computer Graphics Final Project\\[0.3em]
       \large{Real-time Smoke Simulation and Shading}}
\author{蔡平樂 b11201024 \and 許佑威 b11901089}
\date{\today}

\begin{document}
\vspace{-2em}
\maketitle
\vspace{-3em}

% ─────────────────────────────────────────────────────────────────────────────
\section{Introduction}

This project implements a real-time volumetric smoke simulation that runs
entirely in the browser using WebGL.
We start from the open-source 2D fluid simulation by
PavelDoGreat~\cite{pavelfluid} and extend it in several directions:

\begin{enumerate}[leftmargin=*, nosep]
  \item The 2D solver is lifted into a full 3D solver.
  \item A temperature field drives a buoyancy force, and vorticity confinement
        restores small-scale swirling detail lost to numerical dissipation.
  \item The volume is displayed by a per-pixel ray march with soft shadows.
  \item 3D solid geometry—both procedural boxes and meshes loaded from
        GLB files—is integrated into the scene.
  \item A depth pre-pass records the eye-distance of solid surfaces so the ray
        march terminates correctly at opaque boundaries.
  \item A first-person camera and a projectile system let the user explore the
        scene and trigger smoke bursts on impact.
\end{enumerate}

All simulation passes run as GLSL fragment shaders on a full-screen quad,
reading from and writing to floating-point FBOs (framebuffer objects).
The main loop in \texttt{simulation.js} drives the animation via
\texttt{request\-Animation\-Frame}, capping $\Delta t$ at $\tfrac{1}{60}$\,s
to keep the solver stable.

% ─────────────────────────────────────────────────────────────────────────────
\section{Physics Simulation}

We clone the repository~\cite{pavelfluid} as the starting point and rewrite all
shaders to operate on a 3D volume instead of a 2D canvas.

\subsection{3D Volume Stored as a 2D Texture Atlas}

Because WebGL 1.0 does not expose 3D textures, the $64\times64\times64$ grid is
packed into a single $512\times512$ texture by tiling the 64 Z-slices in an
$8\times8$ grid ($64\times8=512$).
Three GLSL helpers in \texttt{shaders.js} are prepended to every shader:
\texttt{encodeUVW} maps a 3D coordinate to an atlas texel,
\texttt{decodeUVW} inverts it, and \texttt{sampleVolume} provides trilinear
interpolation by blending two adjacent Z-slices with \texttt{mix}.

\subsection{Per-Frame Solver Steps}

\texttt{step(dt)} in \texttt{simulation.js} executes the following passes each
frame.

\paragraph{Buoyancy.}
Hot smoke rises and dense smoke sinks (\texttt{passes/buoyancy.js}):
\[
  v_y \mathrel{+}= \Delta t\,\bigl(B\cdot T - W\cdot|\boldsymbol{\rho}|\bigr),
  \quad B=0.1,\quad W=0.001
\]

\paragraph{Vorticity confinement.}
The curl $\boldsymbol{\omega}=\nabla\times\mathbf{v}$ is computed with
six-neighbour finite differences, and a restoring force is injected:
\[
  \mathbf{f} = \varepsilon\,\bigl(\hat{\boldsymbol{\eta}}\times\boldsymbol{\omega}\bigr),
  \quad \hat{\boldsymbol{\eta}} = \operatorname{normalize}(\nabla|\boldsymbol{\omega}|),
  \quad \varepsilon = 30
\]

\paragraph{Pressure projection.}
Incompressibility is enforced in three sub-steps:
(i)~divergence $\nabla\cdot\mathbf{v}$ with six-neighbour differences and
no-slip boundaries;
(ii)~25 Jacobi iterations warm-started at $0.8\times$ the previous pressure:
\[p = \tfrac{1}{6}(p_L+p_R+p_B+p_T+p_F+p_K - \mathrm{div});\]
(iii)~velocity correction $\mathbf{v}\mathrel{-}=\tfrac{1}{2}\nabla p$.

\paragraph{Semi-Lagrangian advection.}
Velocity, density, and temperature are each back-traced and resampled
(\texttt{passes/advection.js}):
\[
  \mathbf{uvw}_{\text{prev}} = \mathbf{uvw} - \Delta t\cdot\mathbf{v}(\mathbf{uvw}),
  \qquad
  q^{n+1} = q^n(\mathbf{uvw}_{\text{prev}})\,/\,(1 + d\,\Delta t)
\]

\subsection{Smoke Emitters}

Smoke is injected by Gaussian splats (\texttt{passes/splat.js}):
$q(\mathbf{uvw}) \mathrel{+}= \mathbf{c}\cdot
\exp(-|\mathbf{uvw}-\mathbf{p}|^2/r)$.
Each emitter is managed by an \texttt{EmitterScheduler}
(\texttt{passes/emitter.js}) supporting a \texttt{startTime}/\texttt{endTime}
window, an optional per-frame \emph{trajectory callback} that overrides
position or intensity, and an initial \emph{burst} (4$\times$ radius,
2$\times$ temperature) that immediately fills the local neighbourhood.

% ─────────────────────────────────────────────────────────────────────────────
\section{Adding 3D Models}

\subsection{GLB File Loading}

\texttt{loadGLBModel} in \texttt{model.js} implements a self-contained binary
glTF~2.0 parser.
It reads the GLB header to locate the JSON chunk and binary buffer chunk,
then reconstructs the scene graph.
Node transforms are decomposed from TRS (translation, quaternion rotation,
scale) into column-major $4\times4$ matrices and multiplied as
$\mathbf{M}=\mathbf{T}\mathbf{R}\mathbf{S}$.
The scene graph is traversed recursively with \texttt{\_traverseNode},
accumulating the world matrix from parent to child.

Accessor data is read via \texttt{\_readAccessor}, handling both
tightly-packed and interleaved buffer views for component types
\texttt{FLOAT32}, \texttt{UINT16}, \texttt{UINT32}, and \texttt{UINT8}.
Material colours are read from the \texttt{pbrMetallicRoughness} base-colour
factor when present; otherwise a default grey-blue $(0.65,0.70,0.80)$ is used.
A bounding box is derived from POSITION accessor \texttt{min}/\texttt{max}
fields to compute a uniform scale that fits a specified target size.

\subsection{Procedural Scene Geometry}

\texttt{createSceneGeometry()} assembles the static scene from box primitives:

\begin{center}
\begin{tabular}{lll}
\toprule
\textbf{Object} & \textbf{Size (W$\times$H$\times$D)} & \textbf{Colour (RGB)} \\
\midrule
Floor            & $16\times0.1\times16$ & Dark grey $(0.18,0.20,0.24)$ \\
Red box          & $2\times2\times2$     & Brick red $(0.84,0.35,0.22)$ \\
Teal box         & $1.5\times1.5\times1.5$ & Teal $(0.28,0.78,0.82)$ \\
Yellow box       & $1\times1\times1$     & Sand $(0.75,0.72,0.36)$ \\
Front/back walls & $8\times4\times0.2$   & Slate $(0.32,0.35,0.38)$ \\
Left/right walls & $0.2\times4\times16$  & Slate $(0.32,0.35,0.38)$ \\
\bottomrule
\end{tabular}
\end{center}

Each box also registers an AABB entry in \texttt{\_colliders} for projectile
collision detection (Section~\ref{sec:input}).

\subsection{GPU Upload and VAO Management}

Every primitive—whether from a GLB or built procedurally—is uploaded by
\texttt{\_buildPrimitive}: separate VBOs for positions and normals, an IBO for
indices, and a VAO (WebGL~2) that encapsulates all attribute state so
subsequent draw calls cannot corrupt each other's bindings.
Index buffers use \texttt{UNSIGNED\_SHORT} or \texttt{UNSIGNED\_INT} depending
on the index array type.
Primitives without normals receive a constant up-vector $(0,1,0)$.

% ─────────────────────────────────────────────────────────────────────────────
\section{Shading}

\subsection{3D Model Shading: Blinn-Phong}

Solid geometry is rendered with the \texttt{\_MODEL\_FRAG} shader in
\texttt{model.js}.
The key-light direction is fixed at $(0.4,0.8,0.45)$ (normalised in-shader)
and the eye position is supplied per-frame via the \texttt{uEye} uniform:
\[
  I = \mathbf{c}_{\text{base}}
      \underbrace{\bigl(\max(\mathbf{N}\cdot\mathbf{L},0)\cdot0.7+0.3\bigr)}_{\text{diffuse + ambient}}
    + \underbrace{\max(\mathbf{N}\cdot\mathbf{H},0)^{48}\cdot0.4}_{\text{specular}}
\]
where $\mathbf{H}=\operatorname{normalize}(\mathbf{L}+\mathbf{V})$.
The constant $0.3$ provides a soft ambient floor; the exponent $48$ gives a
moderately tight specular highlight.
Normals are brought to world space using the upper-left $3\times3$ submatrix
of the model matrix.

\subsection{Smoke Shading: Volumetric Ray Marching}

The smoke volume is rendered by \texttt{drawRayMarch} in \texttt{passes/render.js}.

\paragraph{Camera and ray construction.}
An FPS camera basis $(\hat{F},\hat{R},\hat{U})$ is built from yaw and pitch
angles at runtime:
\[
  \hat{d} = \operatorname{normalize}\!\left(
    \hat{F}
    + \mathrm{ndc}_x\cdot\mathrm{aspect}\cdot\tan\tfrac{\mathrm{FOV}}{2}\cdot\hat{R}
    + \mathrm{ndc}_y\cdot\tan\tfrac{\mathrm{FOV}}{2}\cdot\hat{U}
  \right)
\]
Field of view is $60^\circ$; the aspect ratio is computed from the canvas each
frame.

\paragraph{Ray-box intersection.}
The ray is intersected with the volume AABB (half-extent 5.0, centre
$(0,-2,12)$) using the slab method, yielding $t_{\text{start}}$ and
$t_{\text{end}}$.

\paragraph{Step-and-accumulate loop.}
The interval $[t_{\text{start}},t_{\text{end}}]$ is divided into 64 equal
steps (\texttt{MAX\_STEPS}).
At each sample position $\mathbf{p}$:

\begin{enumerate}[nosep]
\item The volume UVW is $(\mathbf{p}-\mathbf{c}_{\text{vol}})/(2r)+0.5$.
      Samples with density $\rho<0.001$ are skipped.

\item \textbf{Soft shadow.}
      Three steps of length 0.12 are marched toward the light; accumulated
      density attenuates illumination:
      $\ell = \exp(-\sigma_{\text{sh}}\cdot\texttt{DENSITY\_SCALE}\cdot
      \texttt{ABSORPTION}\cdot\texttt{SHADOW\_STEP})$.

\item \textbf{Temperature tint.}
      Sample temperature interpolates between cool blue-white
      $(0.85,0.90,1.00)$ and warm orange $(1.00,0.55,0.10)$:
      $\mathbf{c}_{\text{tint}}=\operatorname{mix}(\text{cool},\text{warm},
      \operatorname{clamp}(T\cdot0.5,0,1))$.

\item \textbf{Front-to-back compositing:}
      \[
        \alpha = 1-e^{-\rho\,\Delta s\cdot\texttt{ABS}},\quad
        \mathbf{c}_{\text{acc}} \mathrel{+}= \tau\,\alpha\,
          \mathbf{c}_{\text{tint}}(\mathbf{c}_{\text{light}}\,\ell+0.35),\quad
        \tau \mathrel{*}= (1-\alpha)
      \]
      Marching stops early when $\tau<0.01$.
\end{enumerate}

The $0.35$ ambient term prevents smoke in full shadow from being completely
invisible.
The result is blended over the background using premultiplied alpha
(\texttt{blendFunc(ONE,\allowbreak\ ONE\_MINUS\_SRC\_ALPHA)}).

% ─────────────────────────────────────────────────────────────────────────────
\section{Depth Buffer Integration}

\subsection{Depth Pre-Pass}

Before ray marching, \texttt{drawModelDepthCapture()} renders all scene
primitives into an off-screen RGBA half-float FBO (\texttt{\_modelScreenFBO})
with an attached 16-bit depth renderbuffer.
The shader (\texttt{\_MODEL\_DEPTH\_FRAG}) writes the eye-to-surface Euclidean
distance into the \emph{alpha} channel:
\[
  \texttt{gl\_FragColor} = \texttt{vec4}(0,\,0,\,0,\,|\mathbf{w}-\mathbf{eye}|)
\]
Alpha zero signals no solid surface at that pixel.
The pass enables \texttt{DEPTH\_TEST} (\texttt{LEQUAL}), depth writes, and
back-face culling, using the same $60^\circ$ perspective and view matrix as
the ray march so the recorded distances are consistent.

\subsection{Ray-March Integration}

Inside \texttt{drawRayMarch}, the pre-pass texture is bound as
\texttt{uModelBuffer}.
For each pixel, the alpha value gives the model depth $d_{\text{model}}$,
which truncates the march interval:
\[
  t_{\text{end}} = \min(t_{\text{AABB}},\; d_{\text{model}})
\]
When alpha is zero, $d_{\text{model}}$ is set to $10^9$ so the ray travels
the full AABB extent.
If $t_{\text{end}}\le t_{\text{start}}$, the fragment outputs transparent
black immediately, saving the cost of the loop entirely.

\subsection{Render Order}

The final frame is assembled in three ordered draw calls:
\begin{enumerate}[nosep]
  \item Depth pre-pass $\to$ off-screen FBO.
  \item Rasterise solid geometry $\to$ default framebuffer (Blinn-Phong,
        depth test on).
  \item Ray march $\to$ default framebuffer (reads pre-pass depth, blended
        over rasterised models with
        \texttt{blendFunc(ONE, ONE\_MINUS\_SRC\_ALPHA)}).
\end{enumerate}

% ─────────────────────────────────────────────────────────────────────────────
\section{Interactive Controls}
\label{sec:input}

\subsection{First-Person Camera}

Camera control in \texttt{input.js} uses the Pointer Lock API for raw mouse
deltas.
Movement basis is derived from the current yaw so that WASD stays in the XZ
plane regardless of pitch.

\begin{center}
\begin{tabular}{ll}
\toprule
\textbf{Key / Input} & \textbf{Action} \\
\midrule
\texttt{W}/\texttt{S}        & Move forward / backward \\
\texttt{A}/\texttt{D}        & Strafe left / right \\
\texttt{Shift}/\texttt{Ctrl} & Move up / down \\
Mouse move (locked)          & Rotate yaw and pitch \\
\texttt{1}--\texttt{4}       & Select smoke-grenade colour (hotbar) \\
Left click                   & Throw the selected smoke grenade \\
\texttt{P}                   & Pause / resume simulation \\
\texttt{Space}               & Toggle depth visualisation \\
\bottomrule
\end{tabular}
\end{center}

Pitch is clamped to $\pm0.44\pi$.
Move speed is 0.05 world units per key event; mouse sensitivity is
0.002\,rad/pixel; both are configurable in \texttt{config.js}.

\subsection{Projectile System: Smoke Grenades}

Left-clicking (after pointer lock) throws a \emph{smoke grenade}—an instance
of a GLB mesh (\texttt{assets/models/ball.glb}) loaded once at start-up by
\texttt{loadProjectileTemplate} and cloned per throw by
\texttt{createProjectileModel}.
Each grenade spawns 0.6 units ahead of the camera with initial velocity
$\mathbf{v}_0 = 8.0\,\hat{F} + (0,1.2,0)$; the upward bias gives it a
natural arc.

\paragraph{Hotbar and grenade types.}
\texttt{PROJECTILE\_TYPES} in \texttt{model.js} defines four grenade colours
(white, red, blue, green), each pairing a model tint with a matching smoke
colour.
Pressing \texttt{1}–\texttt{4} calls \texttt{setProjectileType}, which updates
\texttt{currentProjectileType} and refreshes an on-screen hotbar
(\texttt{updateProjectileHUD}): the selected slot is highlighted and a
\texttt{\#weapon-name} label shows the grenade's display name (e.g.\ ``Red
smoke grenade'').
Cloned primitives reuse the template's geometry buffers and simply override
\texttt{baseColor} with the type's tint, avoiding redundant GPU uploads.

\paragraph{Impact response.}
Every frame, \texttt{updateProjectiles(dt)} integrates gravity
($a_y=-3.0$) and tests the motion segment against each AABB in
\texttt{\_colliders} using the slab method.
On the nearest hit, \texttt{onProjectileHit} converts the impact point to
volume UVW and performs three Gaussian splats (\texttt{splat3D}, radius
0.006):
\begin{enumerate}[nosep]
  \item the grenade's smoke colour into the \emph{density} field, producing a
        coloured smoke burst;
  \item a uniform heat value $(0.35,0.35,0.35)$ into the \emph{temperature}
        field, so the burst rises by buoyancy;
  \item the normalised impact velocity, scaled by $0.4$, into the
        \emph{velocity} field, so the burst inherits the grenade's direction
        of travel as an initial impulse.
\end{enumerate}
Projectiles surviving more than 5 seconds are removed automatically.
Because projectiles participate in the depth pre-pass, they correctly
occlude the smoke volume.

\subsection{Depth Visualisation}

Pressing \texttt{Space} toggles the flag
\texttt{config.SHOW\_DEPTH\_VIZ}.
In this mode, \texttt{drawDepthViz} reads the alpha channel of the pre-pass
FBO and maps eye-distance to greyscale ($1 - d/20$), giving a live view of
the depth buffer for debugging.

% ─────────────────────────────────────────────────────────────────────────────
\section{Limitations and Future Work}

\paragraph{Smoke does not flow around solids.}
The depth pre-pass makes smoke \emph{appear} to stop at surfaces, but the
fluid solver has no solid boundary conditions.
Velocity and density fields pass through geometry unimpeded.
A proper fix would stamp a solid-occupancy mask each frame and enforce
no-slip conditions in the divergence and pressure-gradient passes.

\paragraph{Volume resolution.}
At $64^3$ voxels, fine detail is lost at close range due to trilinear
smearing.
Doubling the linear resolution to $128^3$ would require a $1024\times1024$
atlas (4$\times$ the memory and roughly 8$\times$ the shader work), making
higher resolution expensive without further optimisation such as adaptive
grid refinement.

\paragraph{Shadow quality.}
Only three fixed-length shadow-ray steps are used, which underestimates
occlusion inside thick smoke.
More steps, or an adaptive step size based on local density, would improve
accuracy at a proportional GPU cost.

\paragraph{Single emitter.}
The default scene configures one emitter with a four-second active window.
The \texttt{EmitterScheduler} class already supports multiple simultaneous
emitters with scripted trajectories, so richer smoke sources could be added
with minimal code changes.

% ─────────────────────────────────────────────────────────────────────────────
\section{References}

\begin{thebibliography}{9}
\bibitem{pavelfluid}
  P.\ Dobrev (PavelDoGreat),
  \textit{WebGL Fluid Simulation}.\\
  \url{https://github.com/PavelDoGreat/WebGL-Fluid-Simulation}

\bibitem{firesim}
  A.\ Chan,
  \textit{Real-time Fire Simulation}.\\
  \url{https://andrewkchan.dev/posts/fire.html}
\end{thebibliography}

\end{document}
